\chapter{Conclusion}
\label{ch:conclusion}

Systems of twinned systems represent a novel, emerging class of systems, in which multiple digital twins are composed according to system of systems principles. These systems achieve capabilities beyond those of individual systems. The systematic literature review reported in this thesis demonstrates that systems of twinned systems experience reliability challenges due to dynamic composition and interdependencies between constituents. These challenges, in turn, lead to disturbances that propagate across the system and affect its overall behaviour. In this thesis, I propose a model-driven approach that captures the relationship between constituents' health and preference to belong to the system. My approach relies on Statecharts to codify the overall behaviour of the system. Through this model, I define reliability levels for systems of twinned systems that specify behavioural agreements between constituents and the system. These agreements mitigate disturbances and, through that, improve the reliability of services built on an inherently unreliable system of twinned systems. I demonstrate the utility of the approach through a case of a typical service in distributed systems: complex event processing. The demonstration shows that the proposed approach’s mechanism of controlling when and how constituents contribute, together with compensation, maintains continuous event streams under data loss and improves the reliability of the service. The proposed approach provides a structured way to manage reliability in systems of twinned systems and paves the way for the next generation of complex, distributed cyber-physical ecosystems.